%% CV - Giovanny Garcia
%% Tailored for TBD Investors, Software Engineering Intern
%% Compile with: lualatex main_tbd_investors.tex

\documentclass[11pt,a4paper,sans]{moderncv}
\moderncvstyle{banking}
\moderncvcolor{blue}

% Force first/last name and section headings to moderncv blue (color1).
% These commands exist on newer moderncv (MiKTeX) and are absent on TeX Live
% 2023's moderncv 2.3.1, so only redefine when they are defined.
\makeatletter
\@ifundefined{firstnamestyle}{}{%
  \renewcommand*{\firstnamestyle}[1]{{\fontsize{34}{36}\bfseries\upshape\color{color1}#1}}
  \renewcommand*{\lastnamestyle}[1]{{\fontsize{34}{36}\bfseries\upshape\color{color1}#1}}
  \renewcommand*{\sectionstyle}[1]{{\sectionfont\color{color1}#1}}
}
\makeatother

\usepackage[utf8]{inputenc}
\usepackage{needspace}
\usepackage[scale=0.80]{geometry}
\usepackage{import}

% personal data
\name{Giovanny}{Garcia}
\address{Austin, Texas, United States}{}{}
\email{giovanny.garcia2@g.austincc.edu}
\extrainfo{\href{https://linkedin.com/in/giovanny-garcia-07ab24322}{linkedin.com/in/giovanny-garcia-07ab24322}, \href{https://github.com/giovanny-garcia}{github.com/giovanny-garcia}}

\begin{document}
\hypersetup{
    colorlinks=true,
    linkcolor=blue,
    filecolor=magenta,
    urlcolor=blue,
    pdftitle={Giovanny Garcia - CV},
    pdfpagemode=FullScreen,
}

\makecvtitle

% ============================================================
%     PROFILE STATEMENT
% ============================================================

\vspace{6pt}
\small{Computer Science graduate (Associate of Science, Austin Community College) who builds operating tools, not research demos. I connect REST APIs, persist data in SQLite, and automate repetitive work so other people can run the system without me. I use Claude Code to turn messy admin into a documented workflow. Seeking TBD Investors' paid, remote Software Engineering Intern role (about 20 hours a week) to ship internal tools, integrations, and practical AI across founder-led portfolio companies.}

% ============================================================
%     CORE COMPETENCIES
% ============================================================

\section{Core Competencies}
\vspace{1pt}
\begin{itemize}

\item \textbf{Software development (TypeScript, Node.js):} Ships typed backend features with discord.js, REST clients, structured errors, and local persistence. Secondary coursework skills in C++, C\#, and React.

\item \textbf{APIs, databases, and integrations:} Cache-first REST clients, SQLite schemas (settings, caches, deduplication, quota counters), and a single sync path so daily use does not hammer a third-party system.

\item \textbf{Workflow automation:} Poll loops, reminder windows, and subscribe/unsubscribe flows that turn a manual check into a reliable announcement. Exponential backoff when a pass fails.

\item \textbf{AI-enabled solutions:} Uses Claude Code to evaluate messy problems, draft structured documents, and multiply output on implementation and debugging while keeping the architecture decisions.

\item \textbf{Independent remote delivery:} Owns a project from requirement to documented code. README-level setup, environment-variable contracts, and usage strategy so someone else can run the tool.

\end{itemize}

% ============================================================
%     INDEPENDENT PROJECTS
% ============================================================

\section{Independent Projects}
\vspace{3pt}
\begin{itemize}

\needspace{5\baselineskip}
\item{\cventry{2026}{CS2 Discord Bot (TypeScript, Node.js, SQLite)}{Personal project}{Austin, TX}{}{\vspace{1pt}
\href{https://github.com/giovanny-garcia/hltv-discord-bot}{github.com/giovanny-garcia/hltv-discord-bot}
\begin{itemize}
    \item Built an operating tool that pulls public match data and turns it into slash commands, announcements, and reminders a small group can use without checking the vendor site by hand (Counter-Strike 2 / GGScore API).
    \item Designed a cache-first integration: API calls happen only on an explicit \texttt{/sync}. Slash commands, announcements, and the poll loop read SQLite so the 3-request/day free tier is not wasted.
    \item Modeled the database around real operating needs: guild settings, seen-item deduplication, cached match/results/country datasets, and daily API usage counters (WAL mode, upserts).
    \item Tested failure modes in the poll loop: exponential backoff when a pass fails, typed API errors, and a \texttt{/quota} command so an operator can see remaining calls and cache status.
    \item Documented setup, environment variables, and a free-tier usage plan so another person can run the bot without burning the daily quota on boot. Secrets stay in \texttt{.env} (gitignored); the README says to regenerate a key if one is ever exposed.
\end{itemize}}}

\vspace{3pt}

\needspace{5\baselineskip}
\item{\cventry{2026}{Applied AI job-search workflow (Claude Code)}{Personal project}{Austin, TX}{}{\vspace{1pt}
\href{https://github.com/giovanny-garcia/ai-job-search}{github.com/giovanny-garcia/ai-job-search}
\begin{itemize}
    \item Forked and ran an open-source application framework on Claude Code to evaluate job postings, tailor a CV, and draft cover letters from a structured profile.
    \item Treat Claude Code as an operations tool: summarize source documents, generate structured reports, and turn a messy admin process into a repeatable workflow.
    \item Keep every generated claim tied to a source file so the output stays factual. That is judgment on administrative data, not production ML.
\end{itemize}}}

\end{itemize}

\clearpage
% ============================================================
%     EDUCATION
% ============================================================

\section{Education}
\vspace{1pt}
\begin{itemize}

\needspace{5\baselineskip}
\item{\cventry{2022--2025}{Associate of Science in Computer Science}{Austin Community College}{Austin, TX}{}{\vspace{1pt}
Coursework foundation in programming and software-development fundamentals. LinkedIn-listed skills include C++, C\#, and React. Seeking a paid internship that combines building with real operating businesses.
}}

\end{itemize}

% ============================================================
%     TECHNICAL SKILLS
% ============================================================

\section{Technical Skills}
\vspace{1pt}
\begin{itemize}
\item \textbf{Languages:} TypeScript, JavaScript, SQL (SQLite), C++, C\#.
\item \textbf{Tools and platforms:} Node.js, discord.js, Git/GitHub, SQLite, REST APIs, Claude Code.
\item \textbf{Practices:} quota-aware API clients, caching, structured errors, README-driven setup, iterating from real user constraints.
\item \textbf{Ready to learn on the job:} production web applications, cloud platforms, and data pipelines are not in my public repos. I pick those up on a live problem, which is how the Discord bot and Claude Code workflow got built.
\end{itemize}

% ============================================================
%     ADDITIONAL BACKGROUND
% ============================================================

\section{Additional Background}
\vspace{1pt}
\begin{itemize}
\item \textbf{How I work:} I pick a user problem, constrain the design (quotas, permissions, reminders), then ship a vertical slice and document it. Ready to do that with portfolio-company operators, without waiting for a ticket queue.
\item \textbf{Work setup:} Austin, Texas. US-based. Available about 20 hours a week for the remote internship, or a fuller summer stretch if that helps a portfolio company.
\end{itemize}

% ============================================================
%     LANGUAGES
% ============================================================

\section{Languages}
\vspace{1pt}
\begin{itemize}
\item English.
\end{itemize}

% ============================================================
%     REFERENCES
% ============================================================

\section{References}
\vspace{1pt}
\begin{itemize}
\item{Available upon request.}
\end{itemize}

\end{document}
